% ---
% id: protocol-compiler/def:exact
% kind: def
% title: "Exact admissibility of declared operator data"
% summary: >-
%   Declared binary64 operator data are admitted by sector-map agreement,
%   Hermiticity, positivity, completeness and a completeness sum that rounds
%   nowhere, each decided exactly on the declared numbers. No tolerance is
%   read and no ideal operator is inferred.
% prominence: standard
% status: implemented
% uses: [def:declaration, foundations/def:binary64, foundations/def:exact-admissibility, foundations/lem:exact-accumulation]
% mentions: [thm:typing, def:compile, def:weighted, thm:sifting-and-key, foundations/req:consumer-decides-exactly, foundations/def:declared-order, foundations/rem:summation-order, foundations/req:exact-emission, foundations/req:boundary-admission, req:validator-refusals]
% assumes: []
% refs: []
% provenance: newly-derived
% implements: certiqs_key.foundations.exact:sum_in_declared_order
% lean:
% notes: "The five conditions are exact admissibility of foundations/def:exact-admissibility with constant 1, restated in this module's words; the per-factor reading and the conditions on Alice's data are this module's strengthenings and stay here. The ensemble form is not an instance of the core: it replaces the core's completeness identity by the identity with tau_A, so the family {N_x} is decided under foundations/req:boundary-admission (ii), and no citation of the core covers it."
% ---

Of the five conditions of the definition, four are conditions on an operator
family and one, (S), is a condition on the binary64 additions that form a
completeness sum. Binary64 addition is not associative, so each such sum comes
with the order in which it is accumulated
(\xref{foundations/rem:summation-order}).

Deciding every admission predicate on \tref{declared data} exactly, on the
\tref[declared value]{declared values}, is what \xref{foundations/req:consumer-decides-exactly} requires of
every consumer.
The conditions the definition states on Alice's data, on $\{M_x\}$ and on
$\tau_A$, are producer obligations of this module and not only consumer
checks, since the consuming engine decides them exactly and refuses
otherwise, so a declaration this module emitted without deciding them would
be refused downstream on data this module had passed. For $\{M_x\}$ and
for the trace sum of $\tau_A$ the obligation is
\xref{foundations/req:exact-emission}\,(ii)--(v), and for the Hermiticity,
positivity and unit trace of $\tau_A$ it is
\xref{foundations/req:boundary-admission}\,(ii). Each is decided here, before
emission, and refused with a named reason otherwise
(Requirement~\ref{req:validator-refusals}).

\begin{definition}
\label{def:exact}
Declared operator data are \tref{binary64 numbers}, each the dyadic rational
it denotes and not an approximation of anything
(\xref{foundations/def:binary64}\,(iii)), and every predicate this module
decides on them is decided \emph{exactly}, on those numbers. This definition lists
those predicates, together with the one condition of the list that is a
condition on the arithmetic that forms a sum rather than a predicate on an
operator.

\emph{The declared summation order.} The \tref{compilation map} emits a
total order on every declared alphabet and on $\SecSet$: the
\tref[index map]{label-to-index map} $\idxm$, which numbers the letters of
each declared alphabet, and the sectors, consecutively from $0$
(Definition~\ref{def:compile}, Theorem~\ref{thm:typing}); the
\tref{declared order} of an alphabet is its $\idxm$ order. For a finite
family $\{a_j\}_{j\in J}$ of $n\times n$ matrices indexed by a declared
alphabet $J$ with $\idxm$ order $j_1\prec\dots\prec j_{|J|}$, the
\tref{sum in the declared order} is accumulated as the
\tref{computed sequential sum} (\xref{foundations/def:declared-order})
\begin{equation}\label{eq:declorder}
  \Sigma_0\coloneqq0,\qquad
  \Sigma_i\coloneqq\Sigma_{i-1}+a_{j_i}\ \ (1\le i\le|J|),\qquad
  \Sigma\coloneqq\Sigma_{|J|},
\end{equation}
formed left to right, entrywise, in binary64, with no re-ordering and no
re-association; the trace of a matrix is likewise accumulated over its row
index in increasing order, which is the order of the emitted array.

\emph{The five conditions.} Let $\{a_j\}_{j\in J}$ be a finite family of
block lists against a sector map $\Smap$, one array $a_j^{(s)}$ on
$\Hs{\Qsys_s}$ per $s\in\SecSet$. The family is \tref{exactly admissible}
for $\Smap$ when
\begin{itemize}[leftmargin=2.6em]
\item[(T)] \emph{(sector map)} the sector map the supplied family carries is
  $\Smap$, the map this module holds: the same block kind, the same block
  identifiers in the same order, the same dimensions, so that
  $a_j=\bigoplus_sa_j^{(s)}$ holds \emph{by construction}. This is a type
  condition: it is decided, never measured, and an off-block entry is not
  small but inexpressible;
\item[(S)] \emph{(the sum \tref{rounds nowhere})} writing $\Sigma^{(s)}$ for
  the total $\Sigma$ that \eqref{eq:declorder} produces on the block family
  $\{a_j^{(s)}\}_j$: for each $s$, every one of the $|J|$ entrywise binary64
  additions forming it is exact, so that the computed $\Sigma^{(s)}$ \emph{is}
  the real sum $\sum_ja_j^{(s)}$;
\item[(H)] \emph{(Hermiticity)} $a_j^{(s)}=a_j^{(s)\dagger}$ for every $(j,s)$;
\item[(P)] \emph{(positivity)} $a_j^{(s)}\succeq0$ for every $(j,s)$; the
  matrix being Hermitian by (H), this says that its spectrum is non-negative;
\item[(C)] \emph{(completeness)} $\Sigma^{(s)}=\idon{\Qsys_s}$ for every $s$.
\end{itemize}
Each of (H), (P), (C) is an \emph{identity or an order relation between
exactly represented numbers} and is decided as such: (H) and (C) entry by
entry, (P) by exact symmetric elimination on the block with the zero-pivot
row test (a negative pivot fails, and a zero pivot passes only when the rest
of its row is zero), cubic in $d_s$. Together the five conditions are exact
admissibility for $\Smap$ with constant $1$, that is with every completeness
sum equal to the identity, in the sense of
\xref{foundations/def:exact-admissibility}, the declared order of $J$ being
its $\idxm$ order: (T), (H), (P) and (S) are the conditions of that
definition bearing the same names, and (C), which reads the computed sum, is
its completeness identity, because under (S) the computed sum $\Sigma^{(s)}$
is the exact sum $\sum_ja_j^{(s)}$
(\xref{foundations/lem:exact-accumulation}).

\emph{The same conditions on Alice's data, where there is one block and no
sector map.} $A$ carries no sector decomposition, so a family on $A$ is one
array per index and (T) is empty for it. $\{M_x\}_{x\in\Xset}$ is
\emph{exactly admissible} when (S), (H), (P) and (C) hold of it on $A$ with
$\idon{A}$ on the right of (C); the declared marginal $\tau_A$ satisfies its
own three conditions when $\tau_A=\tau_A^{\dagger}$, $\tau_A\succeq0$ and
$\tr\tau_A=1$ with the trace accumulated in the row order \emph{without
rounding}, the trace being a sum, read below as an identity of real numbers
exactly as a completeness sum is.

\emph{What is not on this list.} A strict inequality on a declared scalar,
such as the condition of Definition~\ref{def:weighted} that every
public-choice weight be positive, is a domain condition on a declared real
number, decided by its sign and not by any of the five. A relation symbol, a
declared tag and the type of a right-hand side are likewise decided and never
measured (D(vi)). And nothing in this definition is a \emph{norm}: no
quantity here is compared against a tolerance window, and no
\tref[defects]{defect}, the entrywise size of the difference a condition
names, is computed.

\emph{The general format.} For a declaration in the general format of
Definition~\ref{def:declaration}\,(DX1)--(DX6) the same conditions are read
as follows. \emph{Per measured factor:} for every $f\in\mathbb{F}$ the
family $\{F^{f}_{v_f}\}_{v_f\in\Vset^{f}}$ is exactly admissible for the
sector map of $f$ --- (T) against that factor's map, and (S), (H), (P), (C)
on that factor's blocks in that factor's declared order. No condition is
decided on a product block \eqref{eq:productfamily}, none being a datum; that
the product family then satisfies (H), (P) and (C) on $\Qsys$ as identities
of real numbers is shown in (DX3). An entry's operator $\hat O_k$ is decided
(T) and (H) against the map of its scope. \emph{Alice's data, block by
block:} with a factor tagged $A$, condition (T) is not empty for Alice's
data: $\tau_A$, the operators of her measurement slot and a general
$\Theta_k$ are block lists against that factor's map. In POVM form
$\{M^{(t)}_x\}_x$ satisfies (S), (H), (P) and (C) on every $A_t$ with
$\idon{A_t}$ on the right of (C); $\tau_A$ satisfies (H) and (P) block by
block and $\sum_t\tr\tau^{(t)}_A=1$, the traces accumulated in the row order
and then over the blocks in their declared order, rounding nowhere.
\emph{Ensemble form:} the family $\{N^{(t)}_x\}_x$ satisfies (S), (H) and
(P) on every $A_t$, and in place of (C) the identity
$\Sigma^{(t)}=\tau^{(t)}_A$ of \eqref{eq:ensembleform}, entry by entry,
between the declared-order sum and the declared block of $\tau_A$; and
$\tau^{(t)}_A\ne0$ for every $t$, a comparison of declared entries with the
datum $0$. These are all the conditions of the form: nothing is decided of
the denoted $\{M_x\}$ of \eqref{eq:ensemblepovm}, its being a POVM a theorem
(Theorem~\ref{thm:sifting-and-key}\,(v)). At one Alice sector and one measured factor in
POVM form these readings are the paragraphs above.
\end{definition}
